\section{Solution}
\label{sec:solution}

The domain-specific language records processes, ports, dependencies, routes,
storage, secrets, health and operational requirements once. It does not name
Kubernetes resources or vendor fields.

The project objective is to define this abstract model and the transformations
that turn it into a complete set of deployment YAML and JSON files.

For each release, CI gives the generator three inputs: the Project Intent
model, per-process environment files and an image lock that maps image names
to pinned digests. It also receives versioned platform data, including naming
and label rules, defaults, ingress and storage classes, the alert catalogue
and cluster capabilities. The generator uses only these inputs. It does not
read live cluster state or access secret values.

A model-to-model transformation turns these inputs into a \emph{Resolved
Deployment}. This complete, non-text model contains exact resource names,
namespaces, paths, image digests, replicas, resources, probes, ports, addresses,
policy peers, routes, secret references and monitoring rules. Renderers
only turn this model into files and make no deployment decisions.

The model-to-text step generates a \emph{Deliverable Set} for each application.
It contains a \texttt{kustomization.yaml}, Kubernetes YAML files and policy JSON
files. After validation, release CI applies the Kubernetes resources with
\texttt{kubectl apply -k}. Platform facts come from one central, versioned
source, so every repository uses the same platform definition.
